\documentclass[authoryear,preprint,12pt]{elsarticle}

%% The amssymb package provides various useful mathematical symbols
\usepackage{amssymb}
%% The amsmath package provides various useful equation environments.
\usepackage{amsmath}
\usepackage{booktabs}

\journal{Social Networks}

\begin{document}

\begin{frontmatter}

%% Title, authors and addresses

%% use the tnoteref command within \title for footnotes;
%% use the tnotetext command for theassociated footnote;
%% use the fnref command within \author or \affiliation for footnotes;
%% use the fntext command for theassociated footnote;
%% use the corref command within \author for corresponding author footnotes;
%% use the cortext command for theassociated footnote;
%% use the ead command for the email address,
%% and the form \ead[url] for the home page:
%% \title{Title\tnoteref{label1}}
%% \tnotetext[label1]{}
%% \author{Name\corref{cor1}\fnref{label2}}
%% \ead{email address}
%% \ead[url]{home page}
%% \fntext[label2]{}
%% \cortext[cor1]{}
%% \affiliation{organization={},
%%             addressline={},
%%             city={},
%%             postcode={},
%%             state={},
%%             country={}}
%% \fntext[label3]{}

\title{Between Cohesion and Division: The Early Modern Venetian Patriciate} %% Article title

%% use optional labels to link authors explicitly to addresses:
%% \author[label1,label2]{}
%% \affiliation[label1]{organization={},
%%             addressline={},
%%             city={},
%%             postcode={},
%%             state={},
%%             country={}}
%%
%% \affiliation[label2]{organization={},
%%             addressline={},
%%             city={},
%%             postcode={},
%%             state={},
%%             country={}}

%%Research highlights
\begin{highlights}
\item We show that the \emph{grandi curti}, the faction that monopolized the dogeship for more than two centuries, enjoyed a structural advantage within the Venetian marriage network.
\item We show that the marriage network remained structurally stable for most of the \emph{grandi curti} monopoly before undergoing a gradual reorganization towards the end of the period, eventually creating structural holes.
\item We show that, as network cohesion declined and structural holes appeared, families occupying high-betweenness positions gained a structural advantage in forming the large electoral coalitions needed to secure the election of one of their members or allies.
\item Contrary to the traditional interpretation in Venetian historiography, we identify the election of Leonardo Donato in 1606, the first doge from a then high-betweenness family, as the earliest indication of the structural transition that preceded the formal end of the \emph{grandi curti} monopoly in 1612. His election marks the beginning of the shift from the dominance of the high-closure \emph{grandi curti} to the growing electoral importance of high-betweenness families.
\end{highlights}


\end{frontmatter}


\end{document}

